Ce troisième livrable nous a permis de dresser un bilan plus approfondi des opérations d'échantillonnage menées en 2025. Nous travaillons actuellement au formatage et à la préparation des données issues de la campagne de 2026, qui vient de se terminer. En 2026, 31 segments routiers ont été suivis, dont 19 n'avaient pas été échantillonnés en 2025. Une fois ces données intégrées, notre jeu de données couvrira donc 49 segments distincts sur les deux années. La plupart des segments suivis en 2026 ont été visités à six reprises, contre deux en 2025, ce qui devrait améliorer la détection des carcasses et accroître le nombre d'observations.

Concernant le suivi de la petite faune, 31 enregistreurs et 30 caméras seront prochainement rapatriés au laboratoire pour l'extraction et l'analyse des fichiers. En 2026, 22 nouveaux milieux humides ont été suivis aux abords des nouveaux tronçons sélectionnés. Ces sites viendront enrichir la base de données de 2025 et pourront être utilisés dans les analyses du livrable final.